% !TEX root = ../../../main.tex
% Block card: preprocessing.window_matrix -- "Sliding windows + features" (Part I, the blocks chapter).
\cardfield{What it does}
Cuts the span into windows of one length and measures each window with a few dozen
single-number summaries, giving a table with one row per window.

\cardfield{When to use it}
Before \blockref{catalogue.cluster} or a classifier, when you want windows compared by what
they are like (how regular, how spiky, how complex) rather than by their exact trace. It is
also how the timescale ladder (1, 10, 60 and 600 minutes) is built for whole recordings,
often on the cluster.

\begin{plainterms}[A window as a row of numbers]
A ten-minute window holds hundreds of samples. Comparing windows sample by sample is fragile.
Instead each window is described by a list of summaries, like describing a person by height,
age and shoe size rather than by a photograph. Two windows with similar summaries are
counted as similar, even if their traces do not line up.
\end{plainterms}

\cardfield{What goes in, what comes out}
In goes a span of one channel. Out comes a set of windows, each with its row of measures:
22 Catch22 statistics, four fast entropies, two slow entropies and, if you ask, scores from
trained models.

\cardfield{Settings that matter}
\emph{window\_min} is the window length in minutes; keep to the ladder (1, 10, 60, 600) so the
matrix can be browsed by timescale. \emph{step\_frac} is the step between windows as a share
of the window; 1.0 means windows touch end to end. \emph{slow\_entropy} is the expensive
stage: it is skipped automatically for windows longer than 4096 samples.

\cardfield{Watch out for}
This block makes no train/test split: if you train on its windows, use
\blockref{preprocessing.sliding_windows} or a pool for the split. Leave \emph{cnn} and
\emph{rf} off for research question 1: those scores come from models trained on the human
labels, so they carry the answer into the features. A window with a missing sample can get a
silently wrong value for some measures (see the appendix entry).
